\documentclass[11pt]{article}
\usepackage{amsmath,amssymb,amsthm,hyperref}
\newtheorem{theorem}{Theorem}
\newtheorem{proposition}[theorem]{Proposition}
\newtheorem{remark}[theorem]{Remark}
\newcommand{\Q}{\mathbb Q}
\newcommand{\CH}{\operatorname{CH}}
\title{GST spatial source selection: rational fibers and spectral monodromy}
\author{Research continuation; hand proofs and uncompiled Lean drafts}
\date{3 November 2026}
\begin{document}\maketitle

\section{Input, target, and the three executed routes}
The input is an arbitrary original smooth projective complex variety $X$,
an integer $0\leq p\leq\dim X$, and a prescribed rational Hodge class
$\alpha\in H^{2p}(X,\Q)\cap H^{p,p}(X)$. The output sought is an actual
$Z_\alpha\in\CH^p(X)_\Q$ with canonical Betti class $\alpha$.
Base revision: \texttt{b403d3aadd6ed79b2822037941e4503ca1788ec3}.
The owner's topology-first mechanism is retained: geometry is realized from
General Space, and the complete causal process must carry actual cycle data.
The three routes are (i) rational repair of the break detector, (ii) an
actual alpha-selected eigencluster source and its algebraic transfer,
and (iii) constructive branching on the full class fiber.
None uses universal Hodge algebraicity as an existence premise.

\section{Rational repair: a scale is not a missing direction}
\begin{theorem}[Saturation and explicit witnesses]
For a $\Q$-linear cycle map $c:N\to H$ and $m\in\Q^\times$,
\[
 m\alpha\in\operatorname{im}c\quad\Longleftrightarrow\quad
 \alpha\in\operatorname{im}c.
\]
If $c(Z)=m\alpha$, the representative is $m^{-1}Z$.
For integer columns $v_i$, an integer target $v$ belongs to their rational
span iff some nonzero positive integer multiple of $v$ belongs to their
integer lattice.
\end{theorem}
\begin{proof}
Apply linearity to $m^{-1}Z$ or $mZ$. For the lattice statement clear the
finitely many denominators of rational coefficients; divide conversely.
\end{proof}
Consequently the previous Round 51 avatar $\mathbb Zu$ versus $2\mathbb Zu$
tests an integral index, not a rational source defect. Its ten odd-multiple
rejections are real integer-lattice results but disappear over $\Q$.
The corrected rational control is $H=\Q^2$, $c(t)=(t,0)$, with detector
$\delta(a,b)=b$. For example $(0,1)$ remains outside the image after every
nonzero rational scaling. This is an abstract control, not a realized
smooth-projective counterexample. The old certificate's shape census is
still construction-guaranteed; its model names do not construct varieties.

\section{An actual selected-cluster source on the regular spectral region}
Use the earlier alpha-dependent construction on a smooth affine Jouanolou
fold $q:J\to X$: rational topological $K$-theory and Stein realization give
holomorphic idempotents encoding a nonzero integral multiple of $q^*\alpha$;
polynomial matrices $a_\alpha$ approximate these on a finite
cohomology-detecting compact core, retaining the local $0/1$ spectral gap.
These precise analytic steps and their scopes are in the October 24
alpha-selected polynomial spectral record. They do not give algebraic
idempotents. Apply the following construction separately to the two matrices
when their character difference is the chosen encoding.

Add an arbitrarily small constant matrix $C$ while preserving the compact
core's contour gap. One can ensure a nonempty regular spectral region:
at one fixed point $u_0\in J$, the discriminant of $a_\alpha(u_0)+C$ is a
nonzero polynomial in $C$ (regular semisimple matrices exist). Its nonvanishing
set is complex dense and meets every sufficiently small ball. For finitely
many connected components use the intersection of these nonempty dense opens.
Let $U\subset J$ be the resulting nonempty open locus with distinct roots.
This does not assert any particular monodromy group. Write the degree as $n$.
Let $W\to U$ be the algebraic finite etale scheme of all ordered roots.
It is an $S_n$-torsor, possibly disconnected, of degree $n!$.
On $W$ the algebraic spectral projectors
\[
 P_i=\prod_{j\ne i}\frac{a_\alpha-\lambda_j I}{\lambda_i-\lambda_j}
\]
have rank one, pairwise orthogonal images $L_i$, and
$\bigoplus_iL_i=\mathcal O_W^n$. Set $x_i=c_1(L_i)$ in actual Chow.
Thus every elementary symmetric polynomial $e_j(x_1,\ldots,x_n)$ is zero
for $j>0$.

For $0\leq r\leq n$, the selector $f_r:T_r\to U$ is the finite etale
quotient $W/(S_r\times S_{n-r})$. It parametrizes subsets of $r$ roots.
The direct sum of their eigenlines descends to an actual rank-$r$ bundle
$F_r\subset\mathcal O_{T_r}^n$, with complementary bundle $Q_r$.
This is genuine algebraic geometry chosen from $a_\alpha$, hence from
the target data; a topological contour by itself is not a global component
of $T_r$.

\begin{theorem}[Every full-selector Chern transfer is neutral]
Let $P$ be a weighted-homogeneous polynomial over $\Q$ of positive
Chow degree in the Chern classes of $F_r$ and $Q_r$. Then
\[
 (f_r)_*P(c(F_r),c(Q_r))=0\quad\hbox{in }\CH^{>0}(U)_\Q.\tag{1}
\]
More generally the same holds for an ordered partition of the eigenvalues
into blocks of any specified ranks, using any positive-degree polynomial
in the Chern classes of the corresponding tautological block bundles and
tracing over the full flag-selector cover.
\end{theorem}
\begin{proof}
Pull (1) back to $W$. Finite etale base change identifies it with the sum
of $P$ over all subsets (respectively all ordered block partitions) of the
root labels. This sum is a symmetric polynomial $S(x)$ of positive degree.
Every symmetric polynomial is a polynomial in $e_1,\ldots,e_n$:
one proves this directly by subtracting the elementary-symmetric monomial
with the same largest lexicographic term, repeating in each finite degree.
The resulting polynomial has no constant term by homogeneity. All $e_j(x)$
vanish, so the pullback is zero. Pullback along $W\to U$ is injective
with rational coefficients: pushing back its pullback multiplies by the
nonzero degree $n!$. This proves the actual Chow assertion, hence its
canonical Betti counterpart. The argument includes the empty or full block;
no division by a Chern class or by a possibly zero target is used.
\end{proof}
This covers products and mixed Chern polynomials on one selector, not only
an additive character of a single eigenline. If the original ambient bundle
is $E$, the same trace is the symmetric polynomial in $c(E)$; it is not
generally zero. Independent matrices with correlated monodromy, nontrivial
ambient bundles, and classes supported on the excluded discriminant have
not been included in (1).

\section{A component is selected by monodromy, not by a local label}
Assume $U$ connected. A connected component $W_0$ of $W$ has monodromy
group $G\leq S_n$; its degree is $|G|$. A connected component of $T_r$
corresponds exactly to a $G$-orbit on $r$-subsets. Its transfer pulls back
to the sum over that orbit, not to one arbitrarily designated subset.
If $G=S_n$, all selectors of fixed ranks are connected and (1) applies
even to the attempt to choose one component.

\begin{theorem}[Alternating monodromy and the orientation threshold]
For $n\geq3$, $A_n$ is transitive on $r$-subsets, and on ordered
partitions having at least one block of size $\geq2$. Thus for $G=A_n$
these cluster/flag selectors are also connected and their transfers vanish
in all positive degrees. For fully ordered eigenlines, arbitrary
$A_n$-invariant polynomials of degree $d<\binom n2$ vanish after imposing
the elementary symmetric relations. The only possible nonconstant
invariant in the rational coinvariant ring is a multiple of
\[
 \Delta(x)=\prod_{i<j}(x_i-x_j),\qquad \deg\Delta=\binom n2.\tag{2}
\]
Its existence as a formal class does not prove a nonzero geometric return.
\end{theorem}
\begin{proof}
For any two partitions choose a permutation taking one to the other.
If it is odd, postcompose by a transposition inside a target block of
size at least two. This preserves the partition and reverses parity.
For proper nonempty subsets $n\geq3$ guarantees such a block; the
empty and whole subsets are singletons already.

For the invariant statement take a transposition $s$ and decompose an
$A_n$-invariant polynomial $F$ into $(F+sF)/2+(F-sF)/2$. The first part
is symmetric. The second changes sign under every odd permutation, so
vanishes on each hyperplane $x_i=x_j$ and is divisible by every
$x_i-x_j$, hence by their product. Its quotient by $\Delta$ is symmetric.
Modulo $(e_1,\ldots,e_n)$ the two symmetric factors reduce to constants.
The same conclusion holds for invariant classes in the quotient: averaging
any representative over the finite group produces an invariant polynomial
representative because the ideal is stable and $|A_n|$ is invertible.
Below the degree in (2) the alternating summand is absent.
\end{proof}
In rank three the cyclic orbit cubic
$C=x^2y+y^2z+z^2x$ satisfies
$C=\Delta/2$ modulo $e_1,e_2,e_3$; the reverse cyclic orbit is
$-\Delta/2$. Their full symmetric trace is zero. This supplies a real
algebraic distinction between orbit selection and full-label averaging.
For $p=2$ and $n\geq3$, even fully ordered Chern sources have zero
positive-degree transfer under $A_n$ in that degree.

\section{A global incidence construction includes the discriminant}
The local transfer admits a second, independent geometric calculation that
does not remove the branch divisor. Let $B=\operatorname{Gr}(r,n)$ and
let $\mathcal F,\mathcal Q$ be its tautological subbundle and quotient.
On $J\times B$, the matrix $a_\alpha$ defines the section
\[
 s_a:\mathcal F\longrightarrow\mathcal Q,\qquad
 s_a(u,F)=a_\alpha(u)|_F\pmod F.
\]
Its zero scheme $Z_a$ is the actual scheme of invariant rank-$r$ subspaces;
$f:Z_a\to J$ is proper. The bundles restrict to genuine algebraic bundles
on $Z_a$. Over $U$ this is exactly the previous finite selector. A separated
analytic spectral cluster supplies a local analytic section of $Z_a$ even
when eigenvalues collide within that cluster.

\begin{theorem}[Global full-incidence neutrality]
One can choose an arbitrarily small constant perturbation $C$ preserving
the compact target gap such that $Z_{a+C}$ is smooth of dimension $\dim J$,
the section is regular, and every irreducible component meets the regular
spectral region. For any such perturbation and every homogeneous polynomial
$P$ of positive Chow degree in $c(\mathcal F),c(\mathcal Q)$,
\[
 f_*P(c(\mathcal F),c(\mathcal Q))=0
       \quad\text{in }\CH^{>0}(J)_\Q,\qquad
 f_*[Z_{a+C}]=\binom nr[J].\tag{4}
\]
For ordered partial flags replace $\binom nr$ by the number of coordinate
flags of the prescribed type; the same positive-degree vanishing holds.
If the regular-region monodromy is transitive on these subsets or flags,
then $Z_{a+C}$ is irreducible, so choosing a whole component does not evade
(4). No claim of that monodromy for every target-dependent matrix is made.
\end{theorem}
\begin{proof}
Write $M=\operatorname{Mat}_n(\C)$. The universal incidence over
$J\times B\times M$ is smooth: at each $(u,F)$, differentiation in $C$
surjects onto $\operatorname{Hom}(F,\C^n/F)$, and its kernel has constant
dimension. The total incidence is therefore an affine-space bundle over
$J\times B$, of dimension $\dim J+n^2$. It dominates $M$ since every
complex matrix has invariant subspaces of every rank. Its generic fiber
is smooth in characteristic zero: localizing the regular coordinate rings
at the generic parameter gives regular rings over the perfect function
field. The nonsmooth locus is closed; its constructible image omits the
generic parameter, so a nonempty parameter open avoids it.

The inverse image of the discriminant in this universal incidence has
dimension $\dim J+n^2-1$. On an incidence fiber over $(u,F)$ it is the
nonzero discriminant polynomial on block-upper-triangular matrices; a
diagonal matrix with distinct entries is a witness to nonvanishing.
The generic-fiber dimension bound, applied to its finitely many components,
gives dimension at most $\dim J-1$ after removing another proper parameter
closed set. Thus no dimension-$\dim J$ component of the smooth $Z_{a+C}$
lies wholly over the discriminant. These two dense parameter opens, and
the regular spectral open at a fixed point, meet every small complex ball.
The compact gap is stable under sufficiently small $C$. The section then
has the expected codimension $d_B=r(n-r)$ and is regular. The same proof
works for partial flags, replacing the displayed Hom bundle by $T_B$:
the infinitesimal $\operatorname{GL}_n$ action surjects on every tangent
space and a regular diagonal matrix has only isolated coordinate flags.

Now the regular zero-section intersection formula gives
\[
 i_*[Z_{a+C}]=\operatorname{pr}_B^*c_{d_B}(T_B)\cap[J\times B],
 \quad i:Z_{a+C}\hookrightarrow J\times B.
\]
Projection formula identifies the push in (4) with the integral over $B$
of $P\,c_{d_B}(T_B)$. For positive degree this is zero since
$\CH^{d_B+\deg P}(B)=0$. For degree zero, a constant diagonal matrix
with distinct eigenvalues has precisely $\binom nr$ reduced zeros on $B$,
so that is the degree of $c_{d_B}(T_B)$. Coordinate flags give the
partial-flag count in exactly the same way.

Finally a smooth variety's irreducible components are disjoint and open.
Each component here meets $U$ and maps dominantly there (the regular-region
map is finite etale onto its connected base). Transitive monodromy makes
the whole selector over $U$ connected, so there can be only one component.
For disconnected original $X$ apply the argument componentwise.
\end{proof}
This is a global vanishing of this \emph{complete} incidence-source recipe,
not a vanishing theorem for every separately selected component, nongeneric
excess cycle, independent mixed-matrix source, or boundary-supported cycle.
For example an algebraic projector with constant eigenvalues may have disconnected invariant-subspace
components whose nonzero individual classes cancel in the total.
For arbitrary nonregular $a$, the top-Chern calculation gives the virtual
zero-locus push, not automatically the fundamental cycle of its excess
components. Those distinctions are essential to the source search.

\section{What this computes for the prescribed original target}
The selector and bundle are actually constructed from $a_\alpha$.
Their canonical Chow/Betti transfer on $U$ is calculated, not postulated.
Full $S_n$ monodromy, or the preceding $A_n$ cases, neutralizes the
regular-region transfer. Equation (4) further neutralizes the complete
generic incidence source globally, including its actual discriminant fibers;
its affine-bundle Chow return to the original $X$ is exactly zero.
For a different extension not governed by that regular incidence formula,
if $\bar Z$ is any Chow extension of
the zero-on-$U$ transfer to $J$, localization says
\[
 \bar Z\in\operatorname{im}\bigl(\CH_{\dim J-p}(J\setminus U)_\Q
       \longrightarrow\CH^p(J)_\Q\bigr).\tag{3}
\]
The discriminant generally has codimension one; (3) therefore gives no
global vanishing and no codimension-two solution. A contour section on a
compact set can meet the repeated-eigenvalue locus internally, provided
its two clusters remain separated. Such a section is not claimed to lie
inside $U$ everywhere. No compact-core class is discarded on that basis.
The remaining geometric routes are independently selected components under
reduced/correlated monodromy, or nongeneric excess/boundary sources not
captured by the complete regular incidence class. The mathematics here
neither constructs them with class
$\alpha$ nor rules them out.

\section{Constructive branching on the full GST fiber}
For an actual CyclePathRealization, form a finite list of rationally
weighted path-transported native seeds and add the resulting cycles.
Linearity and edge naturality prove that its class is the corresponding
sum of transported Hodge vectors. If that sum is a prescribed $\alpha$,
the constructed cycle is its representative; the sum equality is an
explicit checked input, not a theorem of universal coverage.

The distinction between a causal address and its full fiber is geometrically
visible on $Y=\mathbb P^1\times\mathbb P^1$. For its two actual divisor
classes $h_1,h_2$, $h_1^2=h_2^2=0$ and $\int_Yh_1h_2=1$.
The unlabeled ample seed $h=h_1+h_2$ spans only a one-dimensional subspace
in degree two. The nonzero primitive class $h_1-h_2$ has the same
$(2,1,1)$ address and is not a rational multiple of $h$; its square is $-2$.
Its absence from that orbit is not a Hodge counterexample: it is the
class of the actual difference of the two ruling divisors.
The label projectors themselves are genuine algebraic operations:
$L(z)=\pi_1^*(\pi_1)_*(z\cdot h_2)$ and
$R(z)=\pi_2^*(\pi_2)_*(z\cdot h_1)$. They send
$a h_1+b h_2$ to $a h_1$ and $b h_2$, respectively, by the projective
bundle pushforward formula. Thus they are not invented linear lifts.
For any prescribed $a h_1+b h_2$, labeled causal births construct exactly
the cycle $a[\mathrm{pt}\times\mathbb P^1]+b[\mathbb P^1\times\mathrm{pt}]$.
This is a diagnostic refuting universal single-address coverage, not a
replacement of the all-$X$ target by a surface problem.

\section{Proof and evidence boundary}
The selector, descent, Chow transfer, and localization arguments above
are self-reviewed hand proofs. The Lean changes encode rational witnesses,
finite native-cycle branching, rank-three orbit identities, and the
two-ruling multiplicity calculation. They do not instantiate these schemes
or the universal canonical comparison in Lean. Lean compilation is left
to the owner. Exact Wolfram and independent Python computations check
finite invariant/selector identities; the all-rank results follow from the
proofs above, not from extrapolating those tables. No global novelty or
unconditional rational Hodge-conjecture solution is asserted.
\end{document}
